% ---
% id: statistics/lem:exact-strict-tails
% kind: lem
% title: "Exact strict tails at a rational success probability"
% summary: >-
%   At a rational success probability both strict tails of a binomial count are
%   finite sums of binomial probabilities, hence rational numbers, and each
%   comparison of a tail with a rational budget is decided exactly. The endpoint
%   rules can therefore be evaluated without rounding.
% prominence: standard
% status: accepted
% uses: [thm:binomial-tails, cor:tail-forms]
% mentions: [def:conservative-evaluation]
% assumes: []
% refs: []
% provenance: newly-derived
% lean:
% ---

Corollary~\ref{cor:tail-forms} writes the two strict tails as regularised
incomplete beta integrals, the form a numerical evaluation of the incomplete
beta function uses. At a rational success probability the finite sum of the
binomial probabilities gives each tail exactly instead; which of the two
routes an evaluation takes is constrained only by the requirement that it be
a \tref{conservative evaluation}
(Definition~\ref{def:conservative-evaluation}).

\begin{lemma}\label{lem:exact-strict-tails}
Let $n\ge1$, let $p\in[0,1]\cap\mathbb{Q}$ and let $\Ncount{}$ be binomial
with parameters $(n,p)$, as in Theorem~\ref{thm:binomial-tails}. For
$0\le k\le n$,
\begin{equation}\label{eq:tails-sums}
  \Pr[\Ncount{}<k]=\sum_{i=0}^{k-1}\binom{n}{i}p^{i}(1-p)^{n-i},
  \qquad
  \Pr[\Ncount{}>k]=\sum_{i=k+1}^{n}\binom{n}{i}p^{i}(1-p)^{n-i},
\end{equation}
an empty sum being $0$, and:
\begin{enumerate}[label=(\roman*),leftmargin=2.2em]
\item both strict tails are rational numbers, obtained from $n$, $k$ and $p$
by finitely many additions and multiplications of rationals, and they are
equal to the closed forms \eqref{eq:tails-used} of
Corollary~\ref{cor:tail-forms};
\item for every rational $\varepsilon$, whether
$\Pr[\Ncount{}<k]\le\varepsilon$ holds, and whether
$\Pr[\Ncount{}>k]\le\varepsilon$ holds, is decided exactly by evaluating the
rational number \eqref{eq:tails-sums} and comparing two rationals, which is a
comparison of two integers.
\end{enumerate}
\end{lemma}

\begin{proof}
\emph{The sums.} $\Ncount{}$ takes values in $\{0,\dots,n\}$ with
$\Pr[\Ncount{}=i]=\binom{n}{i}p^{i}(1-p)^{n-i}$ (Theorem~\ref{thm:binomial-tails},
with $0^{0}=1$, the convention under which this law is a probability law at
$p\in\{0,1\}$). For $0\le k\le n$ the event $\{\Ncount{}<k\}$ is the disjoint
union of the events $\{\Ncount{}=i\}$ for the integers $0\le i\le k-1$, and
$\{\Ncount{}>k\}$ is the disjoint union of the events $\{\Ncount{}=i\}$ for
$k+1\le i\le n$; at $k=0$, respectively $k=n$, the union is empty and the
event has probability $0$. Finite additivity of the probability gives
\eqref{eq:tails-sums}.

\emph{(i).} Since $p$ is rational, so is $1-p$; every $\binom{n}{i}$ is an
integer; and $p^{i}(1-p)^{n-i}$ is a product of $n$ rational factors (the
empty product being $1$). Each summand of \eqref{eq:tails-sums} is therefore
a product of finitely many rationals, and each tail a sum of at most $n+1$
such summands. The rationals are closed under addition and multiplication, so
both tails are rational. They equal the forms \eqref{eq:tails-used} because
Corollary~\ref{cor:tail-forms} states that the same two probabilities
$\Pr[\Ncount{}<k]$ and $\Pr[\Ncount{}>k]$ equal those forms, for the same
$n$, $p$ and $k$.

\emph{(ii).} By (i) each tail is a rational number $a/b$ with integers $a$
and $b>0$, computed in finitely many exact operations. Write the rational
$\varepsilon$ as $c/d$ with integers $c$ and $d>0$. Then
$a/b\le c/d$ holds if and only if $ad\le cb$, multiplying by $bd>0$, which is
a comparison of two integers and is decided exactly.
\end{proof}
